\documentclass[10pt]{article}
\usepackage[margin=20mm]{geometry}
\usepackage{amsmath,amssymb,mathtools,fontspec,unicode-math,xurl,hyperref}
\IfFontExistsTF{DejaVu Serif}{\setmainfont{DejaVu Serif}}{\setmainfont{Latin Modern Roman}}
\setmathfont{latinmodern-math.otf}
\setlength{\parindent}{0pt}\setlength{\parskip}{6pt}\sloppy
\title{Technical verification of
Sedimentation of particles with very small inertia I}
\author{AI-assisted technical review}
\date{}
\begin{document}
\maketitle
Abstract: This claim-focused review identifies an invalid minimum-distance inference in the proof of the main transport–Stokes limit, an explicit counterexample to the stated unrestricted parameter range of the two general-kernel mean-field theorems, and local errors in two Stokes-field formulas. The first finding is a core proof gap, not a counterexample to the physical limit theorem. The second does not apply to the Oseen kernel. All conclusions concern the acquired arXiv version 5; correspondence with the journal version has not been established. Scientific adjudication remains pending.\par
\section*{Identity, version, and checked scope}
The reviewed work is Richard M. Höfer and Richard Schubert, Sedimentation of particles with very small inertia I: Convergence to the transport-Stokes equation, associated with Duke Mathematical Journal DOI 10.1215/00127094-2024-0038 (https://doi.org/10.1215/00127094-2024-0038). The actual input is arXiv:2302.04637v5, submitted 9 January 2025, whose title page is dated 10 January 2025. It has 59 PDF pages. Throughout this report, page numbers are one-based physical PDF pages, which coincide with the printed numbering.\par
The review traced Theorem 1.1 through the energy estimate, resistance estimates, separation estimate, and bootstrap; examined the assumptions and finite-\(p\) estimates of Section 2; checked Sections 4–5 and the binary and rotational adaptations in Appendices B–C. These are analytical and static-source checks. External results about Stokes well-posedness, solenoidal extension, and convergence of reflections were not reproved from first principles. No author code, particle simulation, or empirical reproduction was run. Detailed coverage and independently written checks accompany this report.\par
The repository page still identifies v5 as its latest version in the check performed for this review. Its note calls v4 an accepted version and describes v5 as correcting an abstract typo. No paper-specific erratum was located in the bounded repository/author-record search. These observations do not establish identity with the unacquired publisher PDF or absence of other corrections.\par
\section*{The mathematical setting and positive checks}
The physical theorem assumes \(\gamma_N=NR\) bounded above and below, \(q>3\), convergence of the initial empirical measure in \(\mathcal W_2\) coupled to the minimum distance as in (H2), the initial relative-velocity condition \[|V_i^0-V_j^0|\le 3\pi\gamma_N\lambda_N |X_i^0-X_j^0|,\] and the energy/maximum-velocity bound (H4). In particular (H2) forces \(\lambda_N\to\infty\). The continuum comparison uses \(\gamma_N\), rather than a fixed limiting value, so omission of a separate \(\gamma_N-\gamma_\ast\) error is intentional.\par
With the unnormalized Euclidean norm, \(|\bar g|=\sqrt N\), the equilibrium velocity is \(\widetilde V=\mathcal R^{-1}\bar g/N\). For \(E=|V-\widetilde V|^2/(2N)\), direct differentiation gives the dissipative coefficient \(\lambda_N N c_{\mathcal R}\). Under the bootstrap estimates \(c_{\mathcal R}\gtrsim N^{-1}\) and \(\|\nabla\mathcal R^{-1}\bar g\|\lesssim N\), Proposition 3.1 yields a normalized velocity error of order \[N^{-1/2}|V-\widetilde V|(0)e^{-c\lambda_Nt}
       +\min\{t,\lambda_N^{-1}\}.\] Its time integral has the order used in (3.37)–(3.41). The extra square on a gradient norm in one intermediate display of the proof is not needed: integrating the preceding differential inequality gives the stated linear dependence directly. It is not treated as an additional substantive problem.\par
The variational resistance scaling \(c_{\mathcal R}\gtrsim R/(1+RS_1)\), \(C_{\mathcal R}\lesssim R\), and the powers in the force-average errors are consistent after the local field corrections below. The finite-\(p\) interpolation exponents in Lemma 2.4 were checked on their stated domain \(0<\beta<d\), \(q>d/(d-\beta)\). These checks do not remove the separation gap.\par
\section*{Findings}
Problem 1 (The minimum-distance bootstrap uses an invalid lower-bound inference).\par
\subsection*{Source claim and exact dependence}
Equation reference: eq:lower. Lemma 3.4, (3.18), p. 26, claims \[
 d_{\min}(t)\ge \tfrac12d_{\min}(0)
 -C\int_0^t d_{\min}(s)A(s)\,ds,
 \quad A=|F|_\infty S_2+R^3S_2S_6^{1/2}|F|_2.\] Its proof ends on p. 44 after (5.16). On p. 29, Step 3 substitutes the force bounds into [eq:lower] (\#eq:lower) and invokes Gronwall to infer (3.44), an \(N\)-uniform exponential lower bound. Step 4, p. 30, explicitly uses (3.44) in (3.45)–(3.46) to close the bootstrap for every fixed time interval. This is a direct dependency of the proof of Theorem 1.1, including its separation conclusion (1.11).\par
\subsection*{Complete check of the suspect steps}
Equation reference: eq:pair. For a fixed pair, Duhamel’s formula and (5.7) give \[
 d_{ij}(t)\ge \tfrac12d_{ij}(0)
       -C\int_0^t d_{ij}(s)A(s)\,ds.\] Choosing a pair that minimizes distance at the final time does not justify replacing its historical distances \(d_{ij}(s)\) by \(d_{\min}(s)\). Since the integral is subtracted, that replacement raises the lower bound. A pair minimizing at time \(t\) need not minimize earlier. The displayed proof supplies no further comparison that establishes this step.\par
Even accepting [eq:lower] (\#eq:lower), a lower integral inequality with a negative integral is not the usual lower differential Gronwall inequality. A complete scalar counterexample has precisely the factor one-half in [eq:lower] (\#eq:lower). Set \[y(t)=2-t^2,\quad b(t)=1,\quad a=\tfrac12y(0)=1,
 \qquad 0\le t\le 7/5.\] The function \(y\) is smooth, positive, and nonincreasing. The margin in \(y(t)\ge a-\int_0^t b(s)y(s)\,ds\) equals \[g(t)=1+2t-t^2-t^3/3.\] Since \(g''(t)=-2-2t<0\), its minimum on the interval occurs at an endpoint; the endpoint values are \(1\) and \(347/375>0\). Nevertheless \[y(7/5)=1/25<e^{-7/5}=0.2465969639\ldots .\] More generally \(y_\varepsilon(t)=2-(2-\varepsilon)t^2\), \(0<\varepsilon<1\), on \([0,1]\) satisfies the same integral premise with the same initial value and \(b=1\), while \(y_\varepsilon(1)=\varepsilon\). The integral premise therefore supplies no uniform positive terminal lower bound, even allowing a different fixed constant in the exponential. The exact polynomial calculations were independently checked in `native\_resistance/independent\_checks.json`.\par
\subsection*{Counterevidence, impact, and repair boundary}
This is a counterexample to the inference, not to the particle–Stokes system. The original relative dynamics contain information lost in [eq:pair] (\#eq:pair). For instance, writing \(r=X_i-X_j\), \(a_N=6\pi\gamma_N\lambda_N\), one has \[r''+a_Nr'=a_N b,\qquad b=(u)_i-(u)_j.\] If \(|b|\le L|r|\) with a uniform \(L\le a_N/4\), then the upper Dini derivative of \(z=|r'|/|r|\) obeys \[D^+z\le z^2-a_Nz+a_NL.\] The cone \(z\le a_N/2\), initially guaranteed by (H3), is invariant, and gives \[|r(t)|\ge |r(0)|\exp[-2z(0)/a_N-2Lt].\] This is a valid conditional repair route. The paper’s displayed coefficient is only bounded by \(C(1+|V^0|_\infty e^{-\gamma_N\lambda_Nt})\). Assumption (H4) bounds \(|V^0|_\infty/\lambda_N\) by an arbitrary finite constant, not the small constant needed for that cone argument. Thus this review has not verified a repair under the original (H1)–(H4).\par
Appendix B, Step 4, p. 49, repeats the pair-minimum and lower-integral inference for the binary system. There a uniform Lipschitz coefficient \(L\) on the intended bootstrap permits the cone argument for \(\lambda_N>4L\), yielding \(d_{\min}(t)\ge d_{\min}(0)e^{-1-2Lt}\). That conditional local repair does not resolve the separate parameter-domain failure below, and cannot simply be imported into the Stokes initial layer.\par
Assessment: Major, a demonstrated invalid derivation in a core proof chain. Theorem 1.1 has not been disproved. A valid separation argument through the initial layer is still required; the unaffected conditional energy and Wasserstein estimates remain useful.\par
Problem 2 (The stated general-kernel domain admits nonconvergent quadratic shear).\par
\subsection*{Claim, assumptions, and source locations}
Theorem 2.1, pp. 13–14, assumes \[|K(x)|+|x||\nabla K(x)|\le C|x|^{-\alpha},\qquad
 \operatorname{div}K=0,\qquad \alpha<d-1,\] and imposes no lower bound on \(\alpha\), no uniform support bound, and no higher-moment condition beyond the selected Wasserstein distance. Under (2.3), it asserts \(\mathcal W_p(\sigma_N(t),\sigma(t))\le e^{Ct}\mathcal W_p(\sigma_N(0),\sigma^0)\). Theorem 2.3, pp. 14–15, inherits this kernel domain and asserts the corresponding vanishing-inertia estimate under (D1)–(D3).\par
\subsection*{An admissible family and its exact evolution}
Take \(d=3,\alpha=-2,p=1,q=2\), and \[K(x_1,x_2,x_3)=(x_2^2,0,0),\qquad
 \sigma^0=\mathbf 1_{[0,1]^3}.\] The kernel is smooth and divergence free, \(K(0)=0\), and \(|K(x)|+|x||\nabla K(x)|\le3|x|^2\). Also \(q=2>3/4=d/(d-\alpha-1)\). The continuum solution is the volume-preserving shear \[(x,y,z)\longmapsto (x+t(y^2-y+1/3),y,z).\] It exists globally and preserves both total mass and its \(L^2\) norm, equal to one.\par
For \(N=n^3\), start from the midpoint \(n^3\) grid in the unit cube. Move one grid point with second coordinate \(y_0=1-1/(2n)\) to height \(L_N=N^{3/4}\), keeping its first and third coordinates unchanged. All other points remain fixed initially. Coupling each cube cell to its midpoint and then moving this one atom gives \[\eta_N(0)=\mathcal W_1(\sigma_N^0,\sigma^0)
 \le \frac{\sqrt3}{2n}+\frac{L_N-y_0}{N}\longrightarrow0 .\] The initial minimum distance is \(1/n=N^{-1/3}\). With \(q'=2\), the exponent in (2.3) is \(4/5\), and its second factor is \(1+d_{\min}(0)N^{1/3}=2\). Thus the full smallness condition holds.\par
Write \(m_{k,N}=N^{-1}\sum_i y_i^k\). Direct summation gives \[m_{1,N}=\tfrac12+\frac{L_N-y_0}{N},\quad
 m_{2,N}=\tfrac13-\frac1{12n^2}+\frac{L_N^2-y_0^2}{N}.\] Every particle keeps its second and third coordinates. Its first-coordinate velocity is the constant \(y_i^2-2m_{1,N}y_i+m_{2,N}\). There are no collisions: particles with distinct \((y,z)\) remain separated in those coordinates, and particles with identical \((y,z)\) have equal velocities and retain their initial first-coordinate separation.\par
Equation reference: eq:shear. The empirical first-coordinate barycenter has velocity \(2(m_{2,N}-m_{1,N}^2)\), whereas the continuum barycenter has velocity \(1/6\). Initial first-coordinate barycenters agree. The Kantorovich dual test function \(x\mapsto x_1\) therefore yields \[
 \mathcal W_1(\sigma_N(t),\sigma(t))
 \ge t\left|2(m_{2,N}-m_{1,N}^2)-\tfrac16\right|
 =2tN^{1/2}(1+o(1))\] for every fixed \(t>0\). This contradicts the claimed bound whose right side tends to zero.\par
To define a sequence for every integer \(N\), use any \(N\) points of the \(\lceil N^{1/3}\rceil^3\) midpoint grid, omitting only the excess \(O(N^{2/3})\) points, and move one retained point in the same way. The omission and weight normalization have \(\mathcal W_1\) cost \(O(N^{-1/3})\) before the outlier is moved. Separation remains comparable to \(N^{-1/3}\), the bounded-coordinate moments change by \(O(N^{-1/3})\), and the same diverging outlier contribution persists. The exact cube subsequence alone already exhibits the contradiction along this complete family.\par
For Theorem 2.3 choose \(V_i^0=0\) and \(\lambda_N=N\). Conditions (D2) and (D3) are immediate; (D1) is \((\eta_N(0)+N^{-1})^{4/5}/2\to0\) on the cube subsequence. The first-coordinate displacement is \[h_N(t)(y_i^2-2m_{1,N}y_i+m_{2,N}),\qquad
 h_N(t)=t-\frac{1-e^{-Nt}}N.\] Hence the barycenter discrepancy is \(2h_N(t)(m_{2,N}-m_{1,N}^2)-t/6\), again diverging as \(2tN^{1/2}\). Collision avoidance remains exact. Both universal conclusions fail on their written domain.\par
\subsection*{Checks, counterevidence, and impact}
The independent rational-moment script computes, for \(N=4096\), an initial-distance upper bound \(0.178891\) and a time-one lower bound \(127.718231\). These finite evaluations corroborate the formulas; the contradiction is the analytic \(N\to\infty\) argument, not a simulation.\par
The Oseen kernel in the physical theorem has \(\alpha=1\), so this example does not refute Theorem 1.1. Interpreting the general-kernel theorems as intended only for \(0\le\alpha<d-1\) excludes the example and agrees with the positive-singularity estimates used in the proof. That is a substantive restriction of the stated domain. The report verifies that it excludes this failure mechanism; it does not certify every endpoint or every part of a revised theorem. Alternatively, treating growing kernels requires additional growth/moment control and a new proof.\par
Assessment: Major for the false general-kernel guarantees, which are presented as an independent methodological contribution. The effect is bounded to their stated generality, separately from Problem 1.\par
Problem 3 (Two explicit auxiliary Stokes fields do not satisfy their defining properties).\par
In Section 4, p. 34, and (4.6), p. 35, the exterior single-sphere field is written as \[(\Phi-R^2\Delta\Phi/6)G,\qquad
 \Phi(x)=\frac{I/|x|+xx^{T}/|x|^3}{8\pi},\] with interior value \(G/(6\pi R)\). Taking \(G=e_1\), the exterior trace at \(x=Re_1\) is \(e_1/(3\pi R)\), whereas changing the minus sign to a plus gives \(e_1/(6\pi R)\). The wrong sign produces a nonzero jump on an open boundary patch, so the piecewise field is not in \(\dot H^1\). The plus sign is already used in (3.6), p. 23. It restores the correct surface-force solution and leaves the decay orders used in the estimates unchanged.\par
In the proof of Lemma 5.1, p. 40, set \[\Psi[W]=\left(\frac{3r}{2|x|}-\frac{r^3}{2|x|^3}\right)
                 \frac{W\times x}{2},\quad
 U=\nabla\times\Psi,\quad
 P_{\rm bad}=r\nabla(|x|^{-1})\cdot W .\] The claimed exterior Stokes equation fails. At the admissible point \(x=2re_1\), with \(r>0\) and \(W=e_1\), exact differentiation gives \[-\Delta U+\nabla P_{\rm bad}=\frac{5}{8r^2}e_1\ne0.\] This is not a constant pressure gauge discrepancy. The compatible pressure is \[P_{\rm good}=\frac{3r(W\cdot x)}{2|x|^3}
            =-\frac{3r}{2}\nabla(|x|^{-1})\cdot W .\] Symbolic differentiation gives zero exterior residual for this corrected pair. It has \(U=W\) on the sphere and the required traction \(-\sigma[U,P_{\rm good}]n=3W/(2r)\). The correction must also be used in the cut-off pressure and weight construction on p. 41. Choosing the cut-off equal to one on \(B_{d/2}\), possible since \(d\ge4r\), places its derivatives in the prescribed annulus. The reciprocal identity then gives the intended normalized weight.\par
Assessment: Minor, supported local construction defects. These errors change equations and boundary traces, so they are not merely cosmetic. The independently defined weak solutions and the scaling estimates survive the checked local corrections. No additional failure of the resistance propositions or of the physical limit is established by these formulas.\par
\section*{Overall dependence and remaining work}
Problem 1 interrupts the supplied route to the main theorem’s separation estimate and bootstrap closure; it does not falsify the theorem. Problem 2 falsifies two subsidiary general-kernel statements on an explicitly admitted range, while leaving the physical kernel outside the counterexample. Problem 3 has verified local corrections and should not be used to inflate the impact of Problems 1–2.\par
The rotational refinement in (1.10) and (1.12) uses \(R^2(\Omega-\widetilde\Omega)\), while the displayed Appendix C energy controls \(R(\Omega-\widetilde\Omega)\). Rotlet scaling offers a plausible extra power of \(R\) in translational and local-\(L^2\) effects. No admissible counterexample to that refinement was established, and a complete sharper block-coupling estimate was not verified. This remains an explicitly unverified aspect, not an additional confirmed error. The finite-\(q\) local-fluid estimates in Step 6 likewise were not completely reconstructed after all interpolations. Neither unresolved item raises severity.\par
Item | Issue judgment | Repair verification | Adoption / remaining work\par
Minimum distance | Major proof gap | Conditional cone only | Unknown; original initial-layer argument needed\par
Kernel range | Major false guarantees | Counterexample exclusion only | Unknown; revised range/proof needed\par
Stokes fields | Minor local defects | Field and boundary correction checked | Unknown; synchronize formulas\par
Scientific decision | Pending | Initial AI review | No human reviews supplied\par
9 R. M. Höfer and R. Schubert, Sedimentation of particles with very small inertia I: Convergence to the transport-Stokes equation, arXiv:2302.04637v5 (2025), https://arxiv.org/abs/2302.04637v5. The 59-page PDF and matching released TeX source were consulted. arXiv repository record for 2302.04637, version history and comments, https://arxiv.org/abs/2302.04637, checked 26 October 2026.\par
\section*{Supporting checks}
The source-manifest SHA256 is `a2186fcdd96333feee04f35a2bd401d3f466b51856eafc87dc090b4b1fd6da74`.\par
This initial-review conversation synthesized and checked the arguments. The model-selection UI observation belongs to the external dispatcher. Backend equivalence was not independently verified. No human scientific review or author adoption is asserted.\par
\end{document}
